\clearpage
\fancyhead[L]{\small MERCURIUS IN SOLE VISUS.}
\hypertarget{native-653-lower}{}
\begin{center}\small Göttingische gelehrte Anzeigen. Stück 169. Seite 1688. 1833 October 21.\\[3pt]\rule{.72\linewidth}{.3pt}\end{center}

\emph{Mercurius in Sole visus, of Overgang van Mercurius over de Zon; den 5 Mei \textup{1832} te Utrecht waargenomen door} G. \textsc{Moll}. Amsterdam 1833. Bei C. G. Sulpke. 43 Seiten. 4.

Wir finden in dieser kleinen Schrift die Beobachtungen des letzten Mercursdurchganges in Utrecht auf der Sternwarte, und in der Wohnung des Herrn \textsc{Van Beek}, wie auch in Leiden; Amsterdam und Nimwegen; ferner eine Zusammenstellung aller Beobachtungen früherer Durchgänge, die in Holland oder durch Holländische Astronomen gemacht sind; endlich die in Utrecht und Leyden angestellten Beobachtungen der Bedeckung des Saturn vom Monde am 8ten Mai 1832. Eine ausführlichere Anzeige würde hier überflüssig sein, da das Wesentliche des Inhalts aus Briefen des Hrn. Prof. \textsc{Moll} an Hrn. Etatsrath \textsc{Schumacher} bereits im 229. Stück der Astronomischen Nachrichten bekannt gemacht ist.
